De forma a estimar diferentes formas de engajamento dos alunos durante as aulas monitoradas, a presente seção descreve o cálculo dos indicadores de engajamento. Esses dados indicam o quão partificipativos foram os alunos.

\subsection{Engajamento por Câmera}
A métrica mede a proporção do tempo de sessão durante a qual os alunos mantiveram sua câmera de vídeo ativa, a nível de permissão no navegador. É definida na Equação~\ref{eq:student_camera_engagement}.

\begin{equation}\label{eq:student_camera_engagement}
E_{c,i} = \frac{T_{c,i}}{T_{s,i}}
\end{equation}

O cálculo é feito para um dado aluno $i$, e é feita a razão entre tempo com câmera ativa ($T_{c,i}$) e o tempo total de sessão ($T_{s,i}$), ambos em minutos.

A métrica de Engajamento por Câmera Médio da Turma corresponde à média aritmética dos engajamentos individuais por câmera. É definida na Equação~\ref{eq:student_camera_engagement_avg}, sendo calculada pela soma das proporções individuais $E_{c,i}$ dividida pelo número total de alunos $n$.

% \lstinputlisting[
%     language=python,
%     caption={Cálculo do Engajamento por Câmera},
%     firstline=8,
%     lastline=66,
%     label={lst:calculate_camera_engagement}
% ]{scripts/metricsCalc/participation.py}

\begin{equation}\label{eq:student_camera_engagement_avg}
\bar{E}c = \frac{1}{n} \sum{i=1}^{n} E_{c,i}
\end{equation}

A partir dos \textit{traces} já processados, estima-se para cada aluno o tempo em que a câmera permaneceu ativa, utilizando os intervalos entre registros consecutivos. A razão entre esse tempo e a duração total da sessão define a proporção individual de engajamento por câmera. A métrica final corresponde à média dessas proporções.

% \textbf{Detecção de câmera ativa}:
% \begin{itemize}
%     \item \texttt{cameraEnabled}: Indica se a câmera está habilitada no sistema
%     \item \texttt{cameraStreaming}: Indica se a câmera está transmitindo vídeo ativamente
%     \item \textbf{Lógica OR}: Considera câmera ativa se qualquer um dos indicadores for verdadeiro
% \end{itemize}


\subsection{Engajamento por Microfone}
Mede a proporção do tempo de sessão em que os alunos mantiveram o microfone ativo, a nível de permissão no navegador. Para um dado aluno $i$, é definido pela razão entre o tempo com microfone ativo ($T_{m,i}$) e o tempo total de sessão ($T_{s,i}$), conforme Equação~\ref{eq:microphone_engagement_student}.

\begin{equation}\label{eq:microphone_engagement_student}
E_{m,i} = \frac{T_{m,i}}{T_{s,i}}
\end{equation}

A partir das proporções individuais, obtém-se o Engajamento por Microfone Médio da Turma, calculado pela média aritmética dos valores $E_{m,i}$ para todos os $n$ alunos, segundo a Equação~\ref{eq:microphone_engagement_avg}.

% \lstinputlisting[
%     language=python,
%     caption={Cálculo do Engajamento por Microfone},
%     firstline=69,
%     lastline=126,
%     label={lst:calculate_microphone_engagement}
% ]{scripts/metricsCalc/participation.py}

\begin{equation}\label{eq:microphone_engagement_avg}
\bar{E}*m = \frac{1}{n} \sum*{i=1}^{n} E_{m,i}
\end{equation}

Com os \textit{traces} previamente tratados, identifica-se o tempo de microfone ativo ao longo da sessão, com base nos intervalos entre registros consecutivos. A proporção individual de engajamento é obtida dividindo-se esse tempo pelo total da sessão, e a média dessas proporções resulta no valor agregado para a turma.

% \textbf{Detecção de microfone ativo}:
% \begin{itemize}
%     \item \texttt{microphoneEnabled}: Indica se o microfone está habilitado no sistema
%     \item \texttt{microphoneStreaming}: Indica se o microfone está transmitindo áudio ativamente
%     \item \textbf{Lógica OR}: Considera microfone ativo se qualquer um dos indicadores for verdadeiro
% \end{itemize}

\subsection{Participação Voluntária Média}
Corresponde à proporção do tempo de sessão durante a qual os alunos demonstram engajamento ativo através de sinais multimodais combinados. Esta métrica considera a presença simultânea na aba da aula juntamente com a ativação de pelo menos um dispositivo de mídia (câmera ou microfone), representando um comprometimento mais substantivo com a atividade de aprendizagem.

Dado um aluno $i$, a participação voluntária é dada pela razão entre o tempo em participação voluntária $T_{p,i}$ e o tempo total de sessão $T_{s,i}$:

\begin{equation}\label{eq:voluntary_participation_student}
P_{v,i} = \frac{T_{p,i}}{T_{s,i}}
\end{equation}

De posse das proporções individuais, obtém-se o valor $P_{v,i}$, que é a média aritmética dos $n$ alunos:

\begin{equation}\label{eq:voluntary_participation_avg}
\bar{P}*v = \frac{1}{n} \sum*{i=1}^{n} P_{v,i}
\end{equation}

% \lstinputlisting[
%     language=python,
%     caption={Cálculo da Participação Voluntária Média},
%     firstline=129,
%     lastline=193,
%     label={lst:calculate_avg_voluntary}
% ]{scripts/metricsCalc/participation.py}

Após a consolidação dos \textit{traces}, cada intervalo entre registros é analisado para verificar se o aluno estava na aba da aula e com algum canal de engajamento multimodal ativo (câmera ou microfone). Os intervalos que atendem a essa condição compõem o tempo de participação voluntária $T_{p,i}$, enquanto todos os intervalos formam o tempo total de sessão $T_{s,i}$. A razão individual $P_{v,i}$ é então calculada, e a média dessas razões define a métrica final da turma.


% \textbf{Condições para participação voluntária}:
% \begin{itemize}
%     \item \textbf{Presença obrigatória}: Aluno deve estar na aba da aula (\texttt{is\_lecture\_tab()})
%     \item \textbf{Engajamento multimodal}: Pelo menos um dispositivo de mídia ativo:
%           \begin{itemize}
%               \item Câmera habilitada (\texttt{cameraEnabled}) \textbf{OU}
%               \item Câmera transmitindo (\texttt{cameraStreaming}) \textbf{OU}
%               \item Microfone habilitado (\texttt{microphoneEnabled}) \textbf{OU}
%               \item Microfone transmitindo (\texttt{microphoneStreaming})
%           \end{itemize}
% \end{itemize}

% \textbf{Etapas do cálculo}:
% \begin{enumerate}
%     \item \textbf{Agrupamento e ordenação}: Organiza \textit{traces} por aluno cronologicamente
%     \item \textbf{Verificação combinada}: Para cada intervalo, verifica condições de presença + mídia
%     \item \textbf{Acúmulo temporal}: Soma tempos que atendem aos critérios combinados
%     \item \textbf{Razão individual}: Calcula proporção de participação ativa para cada aluno
%     \item \textbf{Média agregada}: Computa média das proporções individuais
% \end{enumerate}